\documentclass[10pt,a4paper]{amsart}
\usepackage{amsmath,amssymb,amsthm,graphicx}
% Retained minimal preamble for the Stacks Project chapter varieties.tex.
% Verbatim excerpt of the chapter preamble: Xy-pic setup (native preamble lines 14-18), the native theorem
% declarations with the mathematical macros (native preamble lines 154-217), and the project own no-op
% definitions of the reference/slogan/history display environments (native wrapper.tex lines 1-3) that the
% counted statement uses for its slogan. Copied mechanically; nothing added, removed or redefined.
% For commutative diagrams we use Xy-pic
\usepackage[all]{xy}

% We use 2cell for 2-commutative diagrams.
\xyoption{2cell}
% Theorem environments.
%
\theoremstyle{plain}
\newtheorem{theorem}[subsection]{Theorem}
\newtheorem{proposition}[subsection]{Proposition}
\newtheorem{lemma}[subsection]{Lemma}

\theoremstyle{definition}
\newtheorem{definition}[subsection]{Definition}
\newtheorem{example}[subsection]{Example}
\newtheorem{exercise}[subsection]{Exercise}
\newtheorem{situation}[subsection]{Situation}

\theoremstyle{remark}
\newtheorem{remark}[subsection]{Remark}
\newtheorem{remarks}[subsection]{Remarks}

\numberwithin{equation}{subsection}

% Macros
%
\def\lim{\mathop{\mathrm{lim}}\nolimits}
\def\colim{\mathop{\mathrm{colim}}\nolimits}
\def\Spec{\mathop{\mathrm{Spec}}}
\def\Hom{\mathop{\mathrm{Hom}}\nolimits}
\def\Ext{\mathop{\mathrm{Ext}}\nolimits}
\def\SheafHom{\mathop{\mathcal{H}\!\mathit{om}}\nolimits}
\def\SheafExt{\mathop{\mathcal{E}\!\mathit{xt}}\nolimits}
\def\Sch{\mathit{Sch}}
\def\Mor{\mathop{\mathrm{Mor}}\nolimits}
\def\Ob{\mathop{\mathrm{Ob}}\nolimits}
\def\Sh{\mathop{\mathit{Sh}}\nolimits}
\def\NL{\mathop{N\!L}\nolimits}
\def\CH{\mathop{\mathrm{CH}}\nolimits}
\def\proetale{{pro\text{-}\acute{e}tale}}
\def\etale{{\acute{e}tale}}
\def\QCoh{\mathit{QCoh}}
\def\Ker{\mathop{\mathrm{Ker}}}
\def\Im{\mathop{\mathrm{Im}}}
\def\Coker{\mathop{\mathrm{Coker}}}
\def\Coim{\mathop{\mathrm{Coim}}}

% Boxtimes
%
\DeclareMathSymbol{\boxtimes}{\mathbin}{AMSa}{"02}

%
% Macros for moduli stacks/spaces
%
\def\QCohstack{\mathcal{QC}\!\mathit{oh}}
\def\Cohstack{\mathcal{C}\!\mathit{oh}}
\def\Spacesstack{\mathcal{S}\!\mathit{paces}}
\def\Quotfunctor{\mathrm{Quot}}
\def\Hilbfunctor{\mathrm{Hilb}}
\def\Curvesstack{\mathcal{C}\!\mathit{urves}}
\def\Polarizedstack{\mathcal{P}\!\mathit{olarized}}
\def\Complexesstack{\mathcal{C}\!\mathit{omplexes}}
% \Pic is the operator that assigns to X its picard group, usage \Pic(X)
% \Picardstack_{X/B} denotes the Picard stack of X over B
% \Picardfunctor_{X/B} denotes the Picard functor of X over B
\def\Pic{\mathop{\mathrm{Pic}}\nolimits}
\def\Picardstack{\mathcal{P}\!\mathit{ic}}
\def\Picardfunctor{\mathrm{Pic}}
\def\Deformationcategory{\mathcal{D}\!\mathit{ef}}
\newenvironment{reference}{}{}
\newenvironment{slogan}{}{}
\newenvironment{history}{}{}
\usepackage{cleveref}
\crefname{originalreference}{original reference}{original references}
\Crefname{originalreference}{Original reference}{Original references}
\makeatletter
\newlabel{schemes-lemma-separated-implies-valuative}{{01KZ}{source}{}{sourcepointer-1}{}}
\newlabel{schemes-lemma-separated-implies-valuative@cref}{{[originalreference][0][]01KZ}{[1][0][]source}{}{}{}}
\makeatother
\crefname{originallemma}{original lemma}{original lemmas}
\Crefname{originallemma}{Original lemma}{Original lemmas}
\makeatletter
\newlabel{lemma-glue-separated}{{09N0}{source}{}{sourcepointer-2}{}}
\newlabel{lemma-glue-separated@cref}{{[originallemma][0][]09N0}{[1][0][]source}{}{}{}}
\makeatother
\crefname{originallemma}{original lemma}{original lemmas}
\Crefname{originallemma}{Original lemma}{Original lemmas}
\makeatletter
\newlabel{lemma-glue-a-bunch-of-local-rings}{{09N4}{source}{}{sourcepointer-3}{}}
\newlabel{lemma-glue-a-bunch-of-local-rings@cref}{{[originallemma][0][]09N4}{[1][0][]source}{}{}{}}
\makeatother
\crefname{originallemma}{original lemma}{original lemmas}
\Crefname{originallemma}{Original lemma}{Original lemmas}
\makeatletter
\newlabel{lemma-points-in-affine}{{09NI}{source}{}{sourcepointer-4}{}}
\newlabel{lemma-points-in-affine@cref}{{[originallemma][0][]09NI}{[1][0][]source}{}{}{}}
\makeatother
\title{Finitely many points with one-dimensional Noetherian local rings on an integral separated scheme lie in an affine open}
\author{The Stacks Project Authors}
\date{stacks-project:varieties.tex:local-snapshot-20261011; GNU FDL 1.2 or later}
\begin{document}
\maketitle
\noindent\textit{Editorial note.} All delivered mathematical text is the author's own, copied verbatim from the pinned Stacks Project chapter file varieties.tex; no statement or proof step was rewritten, reordered or removed.\par\medskip
\noindent\textit{Editorial note.} The proof contains a commutative diagram typeset with Xy-pic; the retained chapter preamble supplies the native Xy-pic setup, so the diagram compiles exactly as in the source.\par\medskip
\noindent\textit{Editorial note.} The counted span cites four cross-chapter labels, each recorded in assembly.external\_references with its canonical tag as a navigation pointer.\par\medskip
\noindent\textit{Editorial note.} The proof cites no bibliography item inside the counted span.\par\medskip
\noindent\textit{Allowed standard prerequisites (editorial).}\begin{itemize}
\item The valuative criterion of separatedness and separatedness via tensor products. Stacks Project, Schemes, Lemma schemes-lemma-separated-implies-valuative (tag 090L); Varieties, Lemma lemma-glue-separated (tag 09MW). The first excludes a valuation ring dominating two of the local rings and the second converts that into surjectivity of A\_i tensor A\_j -> K.
\item Glueing local rings along a field and the affine containment criterion. Stacks Project, Varieties, Lemmas lemma-glue-a-bunch-of-local-rings (tag 09N1) and lemma-points-in-affine (tag 09NJ). The first builds the semi-local ring A with Spec(A\_i) as an open covering and the second converts the resulting morphism into an affine open containing the points.
\end{itemize}\medskip
\ifdefined\hypertarget\hypertarget{sourcepointer-1}{}\fi
\noindent\textit{Source reference.} Original reference 01KZ: sources/stacks-chow/tags.txt, source lines 2068--2068. Cross-chapter standard prerequisite cited by the author inside the counted span; it is used as a pointer and its own proof is not part of this item.\par\medskip
\ifdefined\hypertarget\hypertarget{sourcepointer-2}{}\fi
\noindent\textit{Source reference.} Original lemma 09N0: sources/stacks-chow/tags.txt, source lines 12198--12198. Cross-chapter standard prerequisite cited by the author inside the counted span; it is used as a pointer and its own proof is not part of this item.\par\medskip
\ifdefined\hypertarget\hypertarget{sourcepointer-3}{}\fi
\noindent\textit{Source reference.} Original lemma 09N4: sources/stacks-chow/tags.txt, source lines 12202--12202. Cross-chapter standard prerequisite cited by the author inside the counted span; it is used as a pointer and its own proof is not part of this item.\par\medskip
\ifdefined\hypertarget\hypertarget{sourcepointer-4}{}\fi
\noindent\textit{Source reference.} Original lemma 09NI: sources/stacks-chow/tags.txt, source lines 12216--12216. Cross-chapter standard prerequisite cited by the author inside the counted span; it is used as a pointer and its own proof is not part of this item.\par\medskip
\begin{lemma}
\label{lemma-finite-set-codim-1-points-in-affine}
Let $X$ be an integral separated scheme. Let $x_1, \ldots, x_r \in X$
be a finite set of points such that $\mathcal{O}_{X, x_i}$
is Noetherian of dimension $\leq 1$. Then there exists an affine
open subscheme of $X$ containing all of $x_1, \ldots, x_r$.
\end{lemma}

\begin{proof}
Let $K$ be the field of rational functions of $X$.
Set $A_i = \mathcal{O}_{X, x_i}$. Then $A_i \subset K$ and $K$
is the fraction field of $A_i$. Since $X$ is separated, and
$x_i \not = x_j$ there cannot be a valuation ring $\mathcal{O} \subset K$
dominating both $A_i$ and $A_j$. Namely, considering the
diagram
$$
\xymatrix{
\Spec(\mathcal{O}) \ar[r] \ar[d] & \Spec(A_1) \ar[d] \\
\Spec(A_2) \ar[r] & X
}
$$
and applying the valuative criterion of separatedness
(Schemes, Lemma \ref{schemes-lemma-separated-implies-valuative})
we would get $x_i = x_j$. Thus we see by
Lemma \ref{lemma-glue-separated}
that $A_i \otimes A_j \to K$ is surjective for all $i \not = j$.
By Lemma \ref{lemma-glue-a-bunch-of-local-rings}
we see that $A = A_1 \cap \ldots \cap A_r$ is a Noetherian
semi-local ring with exactly $r$ maximal ideals
$\mathfrak m_1, \ldots, \mathfrak m_r$ such that $A_i = A_{\mathfrak m_i}$.
Moreover,
$$
\Spec(A) = \Spec(A_1) \cup \ldots \cup \Spec(A_r)
$$
is an open covering and the intersection of any two pieces of this
covering is $\Spec(K)$. Thus the given morphisms $\Spec(A_i) \to X$
glue to a morphism of schemes
$$
\Spec(A) \longrightarrow X
$$
mapping $\mathfrak m_i$ to $x_i$ and inducing isomorphisms of local rings.
Thus the result follows from Lemma \ref{lemma-points-in-affine}.
\end{proof}

\end{document}
